\onecolumn

\begin{center}

\begingroup


\fontsize{6.5pt}{9pt}\selectfont

% \setlength\LTleft{-0.5in}
% \setlength\LTcapwidth{\textwidth} % default: 4in (rather less than \textwidth...)      
% \setlength\LTleft{0pt}            % default: \fill
% \setlength\LTright{0pt}           % default: \fill  

\begin{longtable}{p{.05\textwidth}p{.08\textwidth}p{.75\textwidth}}
\caption{The examples of research idea generation results from the proposed full ResearchAgent.} 
\vspace{-0.1in}
\label{tab:examples} \\

\hline 
\textbf{Index} & \textbf{Types} & \hspace{0.05in} \textbf{Texts}
\\ \hline 
\endfirsthead

\multicolumn{3}{c}%
{{\bfseries \tablename\ \thetable{} -- Continued from the previous page}} \\
\hline 
\textbf{Index} & \textbf{Types} & \hspace{0.05in} \textbf{Texts}
\\ \hline 
\endhead 

\multicolumn{3}{r}{\textit{\textbf{Continued on the next page}}} \\ \hline
\endfoot

\hline \hline
\endlastfoot

%%%%%%%%%%%%%%%%%%%%%%%%%%%%%% Examples %%%%%%%%%%%%%%%%%%%
\multirow{4}{*}{\textbf{1}} & 
\multicolumn{1}{p{.08\textwidth}}{\textbf{Input}} & 
\makecell{
    \multicolumn{1}{p{.75\textwidth}}{\textbf{Title}: Knowledge-Augmented Language Model Prompting for Zero-Shot Knowledge Graph Question Answering} \\
    \noalign{\vskip 0.25ex}\cdashline{0-0}\noalign{\vskip 0.75ex}
    \multicolumn{1}{p{.75\textwidth}}{\textbf{Abstract}: Large Language Models (LLMs) are capable of performing zero-shot closed-book question answering tasks, based on their internal knowledge stored in parameters during pre-training. However, such internalized knowledge might be insufficient and incorrect, which could lead LLMs to generate factually wrong answers. Furthermore, fine-tuning LLMs to update their knowledge is expensive. To this end, we propose to augment the knowledge directly in the input of LLMs. Specifically, we first retrieve the relevant facts to the input question from the knowledge graph based on semantic similarities between the question and its associated facts. After that, we prepend the retrieved facts to the input question in the form of the prompt, which is then forwarded to LLMs to generate the answer. Our framework, Knowledge-Augmented language model PromptING (KAPING), requires no model training, thus completely zero-shot. We validate the performance of our KAPING framework on the knowledge graph question answering task, that aims to answer the user’s question based on facts over a knowledge graph, on which ours outperforms relevant zero-shot baselines by up to 48\% in average, across multiple LLMs of various sizes.} \\
    \noalign{\vskip 0.25ex}\cdashline{0-0}\noalign{\vskip 0.75ex}
    \multicolumn{1}{p{.75\textwidth}}{\textbf{Entities}: ['Natural language', 'Learning management system', 'Question answering', 'Wikipedia', 'Artificial intelligence', 'Named-entity recognition', 'Quality assurance', 'Long short-term memory', 'English language', 'Google Books', 'Natural-language programming', 'LMS color space', 'GitHub', 'Logic learning machine', 'Machine translation', 'DBpedia', 'Integrated library system', 'Spanish language', 'Neural machine translation', 'ArXiv', 'Artificial general intelligence', 'Natural-language understanding', 'Nonlinear programming', 'Codex', 'LoRa', 'Adhesive', 'Common Crawl', 'ROUGE (metric)', 'Transformer', 'Tomas Mikolov']} \\
} \\
\cline{2-3}
& \multicolumn{1}{p{.08\textwidth}}{\textbf{Problem}} & 
\makecell{
    \multicolumn{1}{p{.75\textwidth}}{Developing a Multimodal Knowledge-Aware Prompting System for Multilingual Zero-Shot Question Answering across Structured and Unstructured Data Sources} 
} \\
\cline{2-3}
& \multicolumn{1}{p{.08\textwidth}}{\textbf{Method}} &
\makecell{
    \multicolumn{1}{p{.75\textwidth}}{The further refined and enhanced method for developing a multimodal knowledge-aware prompting system for multilingual zero-shot question answering (MMKP-ZSQA) incorporates advanced refinements in the following components to address the requirements for innovativeness, clarity, and generalizability:} \\ \\ 
    \multicolumn{1}{p{.75\textwidth}}{1. Interconnected Multimodal Knowledge Orchestration (IMKO): } \\ 
    \multicolumn{1}{p{.75\textwidth}}{- IMKO is an advanced evolution of SMKI, incorporating state-of-the-art language and vision models to handle data verbalization with context-aware disambiguation methods.} \\
    \multicolumn{1}{p{.75\textwidth}}{- Techniques like attention mechanisms are employed to weigh the relevance of different data pieces when integrated, revealing how textual and visual data augment each other in a synthesized representation.} \\ \\
    \multicolumn{1}{p{.75\textwidth}}{2. Linguistically Inclusive Retrieval Engine (LIRE): } \\ 
    \multicolumn{1}{p{.75\textwidth}}{- LIRE extends EMKA with an emphasis on semantic understanding, using transformer-based models trained on diverse datasets, including idiomatic and cultural nuances across languages.} \\
    \multicolumn{1}{p{.75\textwidth}}{- Specific algorithms to handle linguistic phenomena such as code-switching and transliteration are included, enhancing the application to a broader set of languages and dialects.} \\ \\
    \multicolumn{1}{p{.75\textwidth}}{3. Prompt Learning and Optimization Nexus (PLON): } \\ 
    \multicolumn{1}{p{.75\textwidth}}{- Building on AMPL, PLON focuses on creating a library of optimized prompts categorized by linguistic features and data modalities, using Bayesian optimization algorithms.} \\
    \multicolumn{1}{p{.75\textwidth}}{- It includes domain adaptation techniques and a wider array of meta-learning strategies with case studies for high and low-resource languages, and outlines theoretical frameworks for their implementation.} \\ \\ 
    \multicolumn{1}{p{.75\textwidth}}{4. Cross-Modal Integrative Learning System (C-MILS): } \\
    \multicolumn{1}{p{.75\textwidth}}{- C-MILS advances SCAS by detailing the use of multi-head attention across modalities for effective and scalable reasoning, with exemplifications on how each modality can enhance comprehension synergistically.} \\
    \multicolumn{1}{p{.75\textwidth}}{- The component now incorporates a layer of abstraction to distill knowledge into a modality-agnostic format, aiding reasoning and facilitating interpretability across languages and data types.} \\ \\ 
    \multicolumn{1}{p{.75\textwidth}}{5. User-Centric Adaptation and Privacy Framework (UCAPF): } \\
    \multicolumn{1}{p{.75\textwidth}}{- UCAPF enriches ALCUM with a user-centric design, focusing on interaction protocols that describe user engagement cycles, feedback loops, and privacy-preserving active learning paradigms, with practical workflows and pseudocode.} \\
    \multicolumn{1}{p{.75\textwidth}}{- Aligns with GDPR and other privacy frameworks to fortify trust and address the ethical use of data in dynamic learning environments.} \\ \\ 
    \multicolumn{1}{p{.75\textwidth}}{6. Global Evaluation and Reflective Testing System (GERTS): } \\
    \multicolumn{1}{p{.75\textwidth}}{- GERTS expands CEB by presenting a multi-tiered validation and reflection methodology to adjust system components based on a diverse set of metrics, including fairness, bias, interpretability, and computational efficiency.} \\
    \multicolumn{1}{p{.75\textwidth}}{- Offers a structured approach to cross-cultural evaluation, including the use of demographic and regional diversity in forming test cohorts.} \\ 
} \\
\cline{2-3}
& \multicolumn{1}{p{.08\textwidth}}{\textbf{Experiment}} & 
\makecell{
    \multicolumn{1}{p{.75\textwidth}}{The experiment, named "Refined Experiment for Multimodal Knowledge-Aware Prompting System for Multilingual Zero-shot Question Answering (RE-MKP-ZSQA)", aims to methodically develop and validate an advanced AI system. The experiment is streamlined to address feasibility, clarity, and reproducibility concerns while upholding robustness and validity by adhering to the following refined phases:} \\ \\ 
    \multicolumn{1}{p{.75\textwidth}}{1. Detailed System Implementation Plan:} \\
    \multicolumn{1}{p{.75\textwidth}}{- Provide a publicly accessible project roadmap with specific milestones, resource allocation, and timelines.} \\ \\
    \multicolumn{1}{p{.75\textwidth}}{2. Dataset Curation with Clear Guidelines:} \\
    \multicolumn{1}{p{.75\textwidth}}{- Publish precise annotation guidelines with strategies to prevent bias.} \\
    \multicolumn{1}{p{.75\textwidth}}{- Document the dataset assembly process, including source selection and data processing procedures.} \\ \\
    \multicolumn{1}{p{.75\textwidth}}{3. Transparent System Training:} \\
    \multicolumn{1}{p{.75\textwidth}}{- Offer a detailed training protocol with hyperparameters, optimization strategies, and Bayesian optimization processes used in PLON.} \\ \\
    \multicolumn{1}{p{.75\textwidth}}{4. Structured Zero-Shot Evaluation:} \\
    \multicolumn{1}{p{.75\textwidth}}{- Outline evaluation metrics derived from GERTS with benchmark datasets to test zero-shot capabilities.} \\ \\
    \multicolumn{1}{p{.75\textwidth}}{5. Clearer Interdisciplinary Evaluation Protocol:} \\
    \multicolumn{1}{p{.75\textwidth}}{- Specify the composition of the evaluation committee, criteria for assessments, and methods for integrating the feedback.} \\ \\
    \multicolumn{1}{p{.75\textwidth}}{6. Iterative Improvement with Validation Metrics:} \\
    \multicolumn{1}{p{.75\textwidth}}{- Describe statistical methods for reflective assessment and continuous improvement, aligned with multi-tiered GERTS methodology.} \\ \\
    \multicolumn{1}{p{.75\textwidth}}{7. User-Centric Design and Privacy Compliance Evaluation:} \\
    \multicolumn{1}{p{.75\textwidth}}{- Structure user studies with targeted data points to assess usability and cultural adaptability.} \\
    \multicolumn{1}{p{.75\textwidth}}{- Outline privacy compliance protocols to adhere to international standards.} \\ \\
    \multicolumn{1}{p{.75\textwidth}}{8. Detailed Global Scalability Evaluation Method:} \\
    \multicolumn{1}{p{.75\textwidth}}{- Define evaluation metrics for scalability tests and describe diverse infrastructural setups.} \\ \\
    \multicolumn{1}{p{.75\textwidth}}{9. Enhanced Reporting for Reproducibility:} \\
    \multicolumn{1}{p{.75\textwidth}}{- Commit to creating a comprehensive report with precise specifications, configurations, and instructions for replication purposes.} \\
    \multicolumn{1}{p{.75\textwidth}}{- Utilize GitHub for version-controlled deposition of code and datasets, and arXiv for openly accessible experiment protocols and findings.} \\
}
\\

\midrule

\multirow{25}{*}{\textbf{2}} & 
\multicolumn{1}{p{.08\textwidth}}{\textbf{Input}} & 
\makecell{
    \multicolumn{1}{p{.75\textwidth}}{\textbf{Title}: Test-Time Self-Adaptive Small Language Models for Question Answering} \\
    \noalign{\vskip 0.25ex}\cdashline{0-0}\noalign{\vskip 0.75ex}
    \multicolumn{1}{p{.75\textwidth}}{\textbf{Abstract}: Recent instruction-finetuned large language models (LMs) have achieved notable performances in various tasks, such as question-answering (QA). However, despite their ability to memorize a vast amount of general knowledge across diverse tasks, they might be suboptimal on specific tasks due to their limited capacity to transfer and adapt knowledge to target tasks. Moreover, further finetuning LMs with labeled datasets is often infeasible due to their absence, but it is also questionable if we can transfer smaller LMs having limited knowledge only with unlabeled test data. In this work, we show and investigate the capabilities of smaller self-adaptive LMs, only with unlabeled test data. In particular, we first stochastically generate multiple answers, and then ensemble them while filtering out low-quality samples to mitigate noise from inaccurate labels. Our proposed self-adaption strategy demonstrates significant performance improvements on benchmark QA datasets with higher robustness across diverse prompts, enabling LMs to stay stable.} \\
    \noalign{\vskip 0.25ex}\cdashline{0-0}\noalign{\vskip 0.75ex}
    \multicolumn{1}{p{.75\textwidth}}{\textbf{Entities}: ['Codex', 'Natural language', 'English language', 'United States', 'Question answering', 'Natural-language programming', 'GTRI Information and Communications Laboratory', 'Artificial intelligence', 'LoRa', 'Llama', 'Python (programming language)', 'Learning management system', 'Natural language processing', 'Reinforcement learning', 'LMS color space', 'Wikipedia', 'GitHub', 'Natural-language understanding', 'London, Midland and Scottish Railway', 'Integrated library system', 'Language model', 'Chinese language', 'Lumen (unit)', 'Spanish language', 'English Wikipedia', 'Logic learning machine', 'Gradient descent', 'Alternative public offering', 'Technology transfer', 'Dialogue system']} \\
} \\
\cline{2-3}
& \multicolumn{1}{p{.08\textwidth}}{\textbf{Problem}} & 
\makecell{
    \multicolumn{1}{p{.75\textwidth}}{Developing a Scalable, Domain-Adaptive Test-Time Training Protocol for Low-Resource Language QA Using Small Language Models} \\ 
} \\
\cline{2-3}
& \multicolumn{1}{p{.08\textwidth}}{\textbf{Method}} &
\makecell{
    \multicolumn{1}{p{.75\textwidth}}{1. Selection of Scalable Compact Language Models (CLMs): Identify and evaluate existing CLMs suitable for adaptation, emphasizing models with minimal computational requirements.} \\ \\
    \multicolumn{1}{p{.75\textwidth}}{2. Creation of a Multilingual Test-Time Training (TTT) Framework: Develop a TTT protocol that enables CLMs to adapt to new domains and languages during the inference phase, leveraging unsupervised learning techniques and pseudo-label generation.} \\ \\
    \multicolumn{1}{p{.75\textwidth}}{3. Synthetic and Unsupervised Data Generation: Utilize a combination of unsupervised and synthetic data generation methods to produce multilingual QA pairs, employing techniques such as back-translation and context-based question synthesis.} \\ \\
    \multicolumn{1}{p{.75\textwidth}}{4. Domain-Adaptive Mechanisms: Introduce domain-adaptive components, including feature adaptation layers and meta-learning algorithms, which tailor the model's behavior to new contexts and languages at test time.} \\ \\
    \multicolumn{1}{p{.75\textwidth}}{5. Incremental Language Addition and Dominance Assessment: Start with a subset of linguistically diverse, low-resource languages. Evaluate domain adaptability for each language via an iterative process, ensuring models learn to prioritize resource efficiency.} \\ \\
    \multicolumn{1}{p{.75\textwidth}}{6. Model Robustness and Generalization: Perform robustness tuning (RT) to prepare models for unforeseen linguistic variations and conduct thorough evaluations across multiple domains to ensure models can generalize their learning effectively.} \\ \\
    \multicolumn{1}{p{.75\textwidth}}{7. Human-In-The-Loop Evaluation: Conduct evaluations with native speakers and domain experts to validate the relevance and accuracy of the QA outputs, incorporating feedback into the iterative training process.} \\ \\
    \multicolumn{1}{p{.75\textwidth}}{8. Open-Sourcing and Community Collaboration: Make the TTT protocol, trained models, and evaluation benchmarks publicly available for the research community, fostering collaboration and further innovation.} \\ 
} \\
\cline{2-3}
& \multicolumn{1}{p{.08\textwidth}}{\textbf{Experiment}} & 
\makecell{ 
    \multicolumn{1}{p{.75\textwidth}}{1. Selection and Preparation:} \\
    \multicolumn{1}{p{.75\textwidth}}{- Identify potential compact language models (CLMs) suitable for domain adaptation and test-time training, focusing on those with minimal computational requirement and the ability to be fine-tuned or adapted in an unsupervised manner.} \\
    \multicolumn{1}{p{.75\textwidth}}{- Prepare a diverse set of low-resource languages and corresponding text corpora, ensuring linguistic diversity and sociocultural significance. Select benchmark datasets for these languages if available.} \\ \\
    \multicolumn{1}{p{.75\textwidth}}{2. Training and Adaptation Procedure:} \\
    \multicolumn{1}{p{.75\textwidth}}{- Create a Test-Time Training (TTT) framework that allows selected CLMs to adapt to various domains in the selected low-resource languages during the inference phase.} \\
    \multicolumn{1}{p{.75\textwidth}}{- Implement unsupervised learning techniques and pseudo-label generation to produce QA pairs, utilizing back-translation and context-based question synthesis to generate synthetic datasets for languages with limited or no available QA datasets.} \\
    \multicolumn{1}{p{.75\textwidth}}{- Integrate domain-adaptive components and meta-learning algorithms into the CLMs to enable domain-specific adaptations at test time.
} \\ \\
    \multicolumn{1}{p{.75\textwidth}}{3. Iterative Evaluation and Refinement:} \\
    \multicolumn{1}{p{.75\textwidth}}{- Begin adaptation and training with a single low-resource language and gradually add additional languages, monitoring the domain adaptability and model performance metrics after each addition.} \\ 
    \multicolumn{1}{p{.75\textwidth}}{- Perform robustness tuning and cross-domain evaluations for each CLM and language adaptation to ensure generalizability and prevent overfitting.} \\ \\
    \multicolumn{1}{p{.75\textwidth}}{4. Human-In-The-Loop Assessment:} \\
    \multicolumn{1}{p{.75\textwidth}}{- Enlist native speakers and domain experts to evaluate the relevance and accuracy of the model's QA outputs for each language.} \\ 
    \multicolumn{1}{p{.75\textwidth}}{- Incorporate feedback into the iterative training process, refining and re-adapting the models accordingly.} \\ \\
    \multicolumn{1}{p{.75\textwidth}}{5. Open-Sourcing and Community Feedback:} \\
    \multicolumn{1}{p{.75\textwidth}}{- Make the TTT protocol, adaptive CLMs, evaluation benchmarks, and any synthetic datasets publicly available for the research community.} \\ \\
    \multicolumn{1}{p{.75\textwidth}}{6. Experiment Monitoring and Documentation:} \\
    \multicolumn{1}{p{.75\textwidth}}{- Record all the parameters, datasets, model configurations, and evaluation metrics meticulously to ensure robustness and reproducibility.} \\ 
    \multicolumn{1}{p{.75\textwidth}}{- Document any challenges faced, unexpected results, or adaptions made during the experiment for open-sourcing purposes.} \\ \\
    \multicolumn{1}{p{.75\textwidth}}{7. Data Analysis and Reporting:} \\
    \multicolumn{1}{p{.75\textwidth}}{- Analyze the collected performance data quantitatively, using appropriate statistical methods to compare with non-adaptive baselines.} \\
    \multicolumn{1}{p{.75\textwidth}}{- Report qualitative findings from human-in-the-loop evaluations, interpreting the implications for language model performance in low-resource language domains.}
}
\\

\midrule

\multirow{2}{*}{\textbf{3}} & 
\multicolumn{1}{p{.08\textwidth}}{\textbf{Input}} & 
\makecell{
    \multicolumn{1}{p{.75\textwidth}}{\textbf{Title}: Whole-brain annotation and multi-connectome cell typing quantifies circuit stereotypy in Drosophila} \\
    \noalign{\vskip 0.25ex}\cdashline{0-0}\noalign{\vskip 0.75ex}
    \multicolumn{1}{p{.75\textwidth}}{\textbf{Abstract}: The fruit fly Drosophila melanogaster combines surprisingly sophisticated behaviour with a highly tractable nervous system. A large part of the fly’s success as a model organism in modern neuroscience stems from the concentration of collaboratively generated molecular genetic and digital resources. As presented in our FlyWire companion paper1, this now includes the first full brain connectome of an adult animal. Here we report the systematic and hierarchical annotation of this 130,000-neuron connectome including neuronal classes, cell types and developmental units (hemilineages). This enables any researcher to navigate this huge dataset and find systems and neurons of interest, linked to the literature through the Virtual Fly Brain database2. Crucially, this resource includes 4,552 cell types. 3,094 are rigorous consensus validations of cell types previously proposed in the “hemibrain” connectome3. In addition, we propose 1,458 new cell types, arising mostly from the fact that the FlyWire connectome spans the whole brain, whereas the hemibrain derives from a subvolume. Comparison of FlyWire and the hemibrain showed that cell type counts and strong connections were largely stable, but connection weights were surprisingly variable within and across animals. Further analysis defined simple heuristics for connectome interpretation: connections stronger than 10 unitary synapses or providing >1\% of the input to a target cell are highly conserved. Some cell types showed increased variability across connectomes: the most common cell type in the mushroom body, required for learning and memory, is almost twice as numerous in FlyWire as the hemibrain. We find evidence for functional homeostasis through adjustments of the absolute amount of excitatory input while maintaining the excitation-inhibition ratio. Finally, and surprisingly, about one third of the cell types proposed in the hemibrain connectome could not yet be reliably identified in the FlyWire connectome. We therefore suggest that cell types should be defined to be robust to inter-individual variation, namely as groups of cells that are quantitatively more similar to cells in a different brain than to any other cell in the same brain. Joint analysis of the FlyWire and hemibrain connectomes demonstrates the viability and utility of this new definition. Our work defines a consensus cell type atlas for the fly brain and provides both an intellectual framework and open source toolchain for brain-scale comparative connectomics.} \\
    \noalign{\vskip 0.25ex}\cdashline{0-0}\noalign{\vskip 0.75ex}
    \multicolumn{1}{p{.75\textwidth}}{\textbf{Entities}: ['Virtual Fly Brain', 'Central nervous system', 'Transposable element', 'SUMO protein', 'Kenyon cell', 'Romani people', 'Induced stem cells', 'Ventral nerve cord', 'FlyBase', "Parkinson's disease", 'Virtual Network Computing', 'P element', 'Piwi-interacting RNA', 'Drosophila Genetic Reference Panel', 'Bateson–Dobzhansky–Muller model', 'J. B. S. Haldane', 'ATG7', "Haldane's rule", 'Oxford Nanopore Technologies', 'Drosophila mauritiana', 'Germline', 'PINK1', 'Migratory locust', 'CRISPR', 'Heliconius', 'GINS (protein complex)', 'Parkin (ligase)', 'Lepidoptera', 'Illumina, Inc.', 'Drosophila']} \\
} \\
\cline{2-3}
& \multicolumn{1}{p{.08\textwidth}}{\textbf{Problem}} & 
\makecell{
    \multicolumn{1}{p{.75\textwidth}}{Investigating the Functional Implications of Connectome Variability in Drosophila's Learning and Memory Circuits Across Different Environmental and Genetic Contexts} \\ 
} \\
\cline{2-3}
& \multicolumn{1}{p{.08\textwidth}}{\textbf{Method}} &
\makecell{
    \multicolumn{1}{p{.75\textwidth}}{The proposed method involves a multi-tiered approach that integrates connectomics, behavioral assays, genetic manipulation, and computational modeling to investigate the functional implications of connectome variability in Drosophila's learning and memory circuits. The method consists of the following steps:} \\ \\
    \multicolumn{1}{p{.75\textwidth}}{1. Connectome Mapping and Variability Analysis:} \\ 
    \multicolumn{1}{p{.75\textwidth}}{a. Utilize the Virtual Fly Brain database to identify and compare individual connectomes, focusing on the mushroom body.} \\ 
    \multicolumn{1}{p{.75\textwidth}}{b. Quantify the variability in connection weights and cell type counts using statistical methods and machine learning algorithms to identify patterns of variability.} \\ \\
    \multicolumn{1}{p{.75\textwidth}}{2. Behavioral Assays:} \\
    \multicolumn{1}{p{.75\textwidth}}{a. Design a series of learning and memory tasks for Drosophila, such as olfactory conditioning or visual pattern recognition.} \\
    \multicolumn{1}{p{.75\textwidth}}{b. Test groups of flies with known connectome profiles under controlled environmental conditions to establish baseline behavioral data.} \\ \\
    \multicolumn{1}{p{.75\textwidth}}{3. Environmental and Genetic Perturbations:} \\ 
    \multicolumn{1}{p{.75\textwidth}}{a. Expose different groups of flies to varied learning paradigms and sensory inputs to create environmental perturbations.} \\ 
    \multicolumn{1}{p{.75\textwidth}}{b. Use CRISPR-Cas9 technology to introduce targeted mutations in genes like PINK1 or Parkin, creating genetic perturbations.} \\ 
    \multicolumn{1}{p{.75\textwidth}}{c. Assess the impact of these perturbations on connectome structure using high-resolution imaging and reconstruction techniques.} \\ \\
    \multicolumn{1}{p{.75\textwidth}}{4. Transcriptomic and Spatial Analysis:} \\ 
    \multicolumn{1}{p{.75\textwidth}}{a. Apply single-cell RNA sequencing and spatial transcriptomics to profile gene expression changes in response to environmental and genetic perturbations.} \\ 
    \multicolumn{1}{p{.75\textwidth}}{b. Correlate transcriptomic data with connectome changes to identify molecular pathways associated with structural and functional plasticity.} \\ \\
    \multicolumn{1}{p{.75\textwidth}}{5. Computational Modeling and Network Analysis:} \\ 
    \multicolumn{1}{p{.75\textwidth}}{a. Develop computational models to simulate the effects of connectome variability on neural circuit function.} \\ 
    \multicolumn{1}{p{.75\textwidth}}{b. Use network analysis tools to explore information flow and circuit dynamics, integrating data from related papers on olfactory projection neurons and information flow in the olfactory system.} \\ \\
    \multicolumn{1}{p{.75\textwidth}}{6. Synthesis and Validation:} \\ 
    \multicolumn{1}{p{.75\textwidth}}{a. Integrate findings from behavioral assays, connectome mapping, transcriptomic analysis, and computational modeling.} \\ 
    \multicolumn{1}{p{.75\textwidth}}{b. Validate the proposed models and hypotheses through iterative experimentation and refinement.} \\ 
} \\
\cline{2-3}
& \multicolumn{1}{p{.08\textwidth}}{\textbf{Experiment}} & 
\makecell{ 
    \multicolumn{1}{p{.75\textwidth}}{1. Pre-experimental Setup:} \\
    \multicolumn{1}{p{.75\textwidth}}{a. Establish a Drosophila melanogaster breeding program to ensure a consistent supply of genetically similar flies for the experiments.} \\
    \multicolumn{1}{p{.75\textwidth}}{b. Develop a standardized protocol for rearing flies under controlled conditions to minimize pre-experimental variability.} \\ 
    \multicolumn{1}{p{.75\textwidth}}{c. Select and validate CRISPR-Cas9 constructs for targeted gene editing of PINK1, Parkin, and other genes of interest.} \\ 
    \multicolumn{1}{p{.75\textwidth}}{d. Train personnel in the use of the Virtual Fly Brain database and relevant computational tools for connectome analysis.} \\ \\
    \multicolumn{1}{p{.75\textwidth}}{2. Connectome Mapping and Variability Analysis:} \\
    \multicolumn{1}{p{.75\textwidth}}{a. Randomly assign individual flies to either a control group or various treatment groups (environmental and genetic perturbations).} \\
    \multicolumn{1}{p{.75\textwidth}}{b. Utilize high-resolution imaging techniques to map the connectomes of flies from each group, with a focus on the mushroom body.} \\
    \multicolumn{1}{p{.75\textwidth}}{c. Apply statistical and machine learning algorithms to quantify and compare the variability in connection weights and cell type counts across groups.} \\ \\
    \multicolumn{1}{p{.75\textwidth}}{3. Behavioral Assays:} \\
    \multicolumn{1}{p{.75\textwidth}}{a. Design and validate a series of learning and memory tasks, such as olfactory conditioning and visual pattern recognition, ensuring tasks are sensitive to subtle differences in performance.} \\ 
    \multicolumn{1}{p{.75\textwidth}}{b. Test flies from each group in the behavioral tasks and record performance metrics.} \\ 
    \multicolumn{1}{p{.75\textwidth}}{c. Analyze behavioral data to establish correlations with connectome profiles.} \\ \\
    \multicolumn{1}{p{.75\textwidth}}{4. Environmental and Genetic Perturbations:} \\
    \multicolumn{1}{p{.75\textwidth}}{a. Expose flies to different learning paradigms and sensory inputs to induce environmental perturbations.} \\ 
    \multicolumn{1}{p{.75\textwidth}}{b. Perform gene editing using CRISPR-Cas9 to create genetic perturbations in the treatment groups.} \\ 
    \multicolumn{1}{p{.75\textwidth}}{c. Re-map connectomes post-perturbation to assess structural changes.} \\ \\
    \multicolumn{1}{p{.75\textwidth}}{5. Transcriptomic and Spatial Analysis:} \\
    \multicolumn{1}{p{.75\textwidth}}{a. Collect brain tissue from flies post-behavioral assays and perform single-cell RNA sequencing and spatial transcriptomics.} \\ 
    \multicolumn{1}{p{.75\textwidth}}{b. Analyze transcriptomic data to identify gene expression changes and correlate these with observed connectome and behavioral variations.} \\ \\
    \multicolumn{1}{p{.75\textwidth}}{6. Computational Modeling and Network Analysis:} \\
    \multicolumn{1}{p{.75\textwidth}}{a. Develop computational models to simulate the impact of observed connectome variability on neural circuit function.} \\ 
    \multicolumn{1}{p{.75\textwidth}}{b. Use network analysis to integrate behavioral, connectomic, and transcriptomic data, focusing on information flow and circuit dynamics.} \\ \\
    \multicolumn{1}{p{.75\textwidth}}{7. Synthesis and Validation:} \\
    \multicolumn{1}{p{.75\textwidth}}{a. Integrate findings across all experimental components to formulate a cohesive understanding of the functional implications of connectome variability.} \\
    \multicolumn{1}{p{.75\textwidth}}{b. Validate models and refine hypotheses through additional targeted experiments, informed by initial findings.}
}
\\

%%%%%%%%%%%%%%%%%%%%%%%%%%%%%%%%%%%%%%%%%%%%%%%%%%%%%%%%%%


\end{longtable}


\endgroup

\end{center}

\twocolumn
